\documentstyle[%


righttag%


,11pt%


,amssymb%


,russian%               remove in Eng.


,izvru%                 Set izven% in Eng.


]{amsart}





%





\newcommand{\const}{\operatorname{const}}


\newcommand{\sh}{\operatorname{sh}}


\newcommand{\vn}{$V_n$}


\newcommand{\vnn}{$\ov V_n$}


\newcommand{\vnb}{$V_n(B)$}


\newcommand{\mr}{\leq}


\newcommand{\nr}{\ne}


\newcommand{\rang}{\operatorname{rang}}


\renewcommand{\a}{\alpha}


\renewcommand{\b}{\beta}


\newcommand{\tts}[1]{}





%\hoffset=-15mm%                  For Eng.


\setcounter{page}{36}


\god{2003}


\nomer{11~(498)} %                        Remove in Eng.


%\nomer{Vol.\,47, No.\,11}%               Set in Eng.


%\str{\pageref{begin}--\pageref{end}}%  Set in Eng.


\udk{514.764}





\author{\it В.А.\,Киосак, Й.\,Микеш}


% NB:   ^^ remove in Eng.


\title{О геодезических отображениях пространств Эйнштейна}


\thanks{Работа выполнена при финансовой поддержке гранта Чешской


республики \No\,201/02/0616.}





\date{}


%\pagestyle{myheadings}%         Set in Eng.


\pagestyle{plain} %              Remove in Eng.


\begin{document}


%\thispagestyle{plain}%          Set in Eng.


\maketitle


\label{begin}








\section{Введение}





Диффеоморфизм $f$ риманова пространства \vn\ на риманово


пространство \vnn\ называют {\it геодезическим отображением},


если $f$ отображает любую геодезическую линию \vn\ в


геодезическую линию \vnn. 


Этими отображениями занималось много авторов, см. например,


[1]--[3]. Исследования по геодезическим отображениям последних


лет отражены в обзорной статье [4]. 





Отметим, что на сигнатуру метрик римановых пространств \vn\ не


налагаем ограничения, как принято, например, в [2], [3]. 


Исследования ведутся локально в классе достаточно гладких функций.





{\it Пространства Эйнштейна}, которые характеризуются условиями на


тензор Риччи


$$


R_{ij}=\frac Rn g_{ij},


$$


где $R$ --- скалярная кривизна, $g_{ij}$ -- метрический тензор,


имеют большое значение как в римановой геометрии, так и в


ее приложениях [1], [2],$\dotsc$


Вопросами о геодезическом отображении пространств Эйнштейна 


занималось много геометров (см., напр., 


[2], [4]--[13]).





Напомним итоговый результат А.З.\,Петрова и В.И.\,Голикова [2] о


геодезических отображениях четырехмерных пространств Эйнштейна:


{\it четырехмерные


пространства Эйнштейна $V_4$ непостоянной кривизны с


сигнатурой Минковского  не допускают нетривиальные геодезические


отображения на римановы пространства $\overline V_4$ с


сигнатурой Минковского}. 


Нами доказана теорема, которая обобщает этот результат.





Распространяя методы исследований


геодезических отображений четырехмерных пространств


Эйнштейна сигнатуры Минковского на эйнштейновы пространства


более высоких размерностей $n>4$,


А.З.\,Петров ([2], сс.\,355, 461)  высказал гипотезу:


{\it пространства Эйнштейна $V_n$ $(n>4)$ сигнатуры Минковского,


отличные от пространств постоянной кривизны, не допускают


нетривиальных геодезических отображений на пространства


Эйнштейна той же сигнатуры.}


Приводим пример, который эту гипотезу опровергает.








Как оказалось, пространства Эйнштейна, допускающие геодезические


отображения, являются по необходимости пространствами $V_n(B)$


[4], [11]--[13].  Указанные


пространства обобщают введенные в [14]


пространства $V(K)$. Поэтому предварительно приводим новые


результаты в теории геодезических отображений пространств


$V_n(B)$, которые затем применяем для изучения геодезических


отображений пространств Эйнштейна.








\section{О римановых пространствах $V_n(B)$}





Риманово пространство \vn, допускающее нетривиальное


геодезическое отображение, будем обозначать через $V_n(B)$


[4], [11], [13], [15], если в нем выполняются уравнения


\begin{equation}


\begin{split}


\text{(а)} \ a_{ij,k} &= \lambda_{i\,}g_{jk}+\lambda_{j}g_{ik},\\


\text{(б)} \ \ \lambda_{i,j} &= \mu g_{ij}+B a_{ij},


\end{split}


\end{equation}


где ``,'' обозначает ковариантную производную в \vn \ 


относительно невырожденного симметрического тензора $a_{ij}$,


ковектора $\lambda_{i}$ и инвариантов  $\mu$ и $B$.








В [11], [15] доказано, что если \vnb\ допускает


любое другое геодезическое отображение на некоторое \vnn, то при


этом отображении также будут выполняться уравнения (1)


при одинаковом инварианте $B$.


Там же была установлена замкнутость пространств $V_n(B)$


относительно геодезических отображений, т.\,е. доказано, что


если $V_n(B)$ допускает геодезическое отображение на некоторое 


$\overline{V}_n$, то $\overline{V}_n$ является пространством 


$\overline{V}_n(\overline B)$.





Установлено [4], что эйнштейновы, обобщенно


полусимметрические и обобщенно Риччи полусимметрические $V_n$,


допускающие нетривиальные геодезические отображения, являются


пространствами $V_n(B)$. Пространствами $V_n(B)$ будут также


пространства $V(K)$ А.С.\,Солодовникова [14] и


Г.И.\,Кручковича [16].





В [3] введено понятие {\it степени подвижности}


относительно геодезических отображений. Доказано


[17], что римановы пространства $V_n$, имеющие степень подвижности


относительно геодезических отображений больше


двух, являются пространствами $V_n(B)$,  $B$ --- $\const$.





Отметим несколько свойств пространств \vnb.





\begin{thm}


Любое геодезическое отображение риманова пространства \vnb, 


$B\neq 0$, является либо нетривиальным, либо гомотетическим.


\end{thm}





\begin{pf}


Предположим, что пространство \vnb,


$B\neq0$, допускает аффинное (т.\,е. тривиальное геодезическое) 


отображение. Тогда существует решение уравнений (1) при 


$\lambda_i=0$.


Из (1(б)) следует $\mu g_{ij}+Ba_{ij}=0$.


При $B\neq 0$ вытекает, что $a_{ij}=-\frac{\mu}{B} g_{ij}$. 


Отсюда согласно ([3], сc.\,75, 121) следует, что геодезическое


отображение является гомотетическим. 


\end{pf}





\begin{thm}


Если в римановом пространстве $V_n(B)$  среди


векторов $\lambda_i$, удовлетворяющих $(1)$, есть ненулевой


изотропный вектор, то $B=0$.


\end{thm}





\begin{pf}


Пусть в \vnb, $B\neq0$, тензор $a_{ij}$,


ковектор $\lambda_i$ $(\neq0)$ и инвариант $\mu$ являются


решением системы уравнений геодезических отображений, которые в \vnb\


имеют вид (1). 





Далее предполагаем, что $\lambda_i$ --- изотропный


вектор, т.\,е. выполняются условия


\begin{equation}


\lambda_\alpha \lambda_\beta g^{\alpha \beta }=0.


\end{equation}


Здесь и дальше $g^{ij}$ являются компонентами обратной матрицы к


$\| g_{ij}\|$.





а) Сначала рассмотрим случай, когда $B\neq\const$.


Из [18] вытекает, что тензор


$a_{ij}$ и вектор $\lambda_i$ удовлетворяют условиям


\begin{equation}


\text{(а)} \ a_{ij} = \a g_{ij}+\b \lambda_i\lambda_j,\quad


\text{(б)} \ \lambda_{i,j} = \gamma g_{ij}+\delta\lambda_i\lambda_j,


\end{equation}


где $\a$, $\b$, $\gamma$, $\delta$ --- некоторые функции инварианта


$\lambda$. 





Дифференцируя (2), подставляя (3(б)), убедимся, что


$\gamma=0$. Тогда, дифференцируя (3(а)), после подстановки


(1) и (3) имеем


$\lambda_i g_{jk}+\lambda_j g_{ik}=\lambda_k


(\a'g_{ij}+(\b'+2\delta)\lambda_i\lambda_j)$.


Отсюда вытекает, что $\lambda_i=0$. В противном случае 


$\rang\| g_{ij}\|\mr2$. 


Следовательно, случай а) доказан.





б) Осталось рассмотреть случай, когда


$B\equiv\const\neq0$. 


Тогда имеют место уравнения (1) и для инварианта $\mu$


выполняется условие


\begin{equation}


\mu_{,i}=2B\lambda_i.


\end{equation}





Дифференцируем (2). После подстановки (1(б)) получаем


\begin{equation}


\mu\lambda_i+B a_{i\a}\lambda^\a=0. 


\end{equation}


Затем дифференцируем (5) и, учитывая (1), (2) и (4), имеем


\begin{equation}


3B\lambda_i\lambda_j+\mu^2g_{ij}+2\mu B a_{ij}+B^2 a_{i\a}a^\a_j=0.


\end{equation}


Продифференцируем (6) по $x^k$, на основании (1) и (4) будем


иметь 


\begin{equation}


4\lambda_k(B\mu\,g_{ij}+ B^2 a_{ij})+\lambda_i c_{jk}+


\lambda_j c_{ik} =0,


\end{equation}


где $c_{ij}$ --- некоторый симметрический тензор.


Так как $\lambda_{i}\neq0$, то существует $\varepsilon^i$ такой,


что $\varepsilon^i\lambda_i=1$. Свертывая (7) 


с $\varepsilon^i\varepsilon^j$, убедимся, что


$c_{\a k}\varepsilon^\a=c \lambda_k$, где $c$ --- некоторый


инвариант. 


Тогда после свертывания (7) с $\varepsilon^i$ получим


$c_{jk}=c_j\lambda_k$, где $c_j$ --- некоторый вектор.


Тензор $c_{jk}$ является симметричным тензором, поэтому имеем


$c_{jk}=\a\lambda_j\lambda_k$, где $\a$ --- некоторый


инвариант. 


Но при этом из (7) вытекает справедливость формулы


\begin{equation}


\mu B\,g_{ij}+ B^2 a_{ij}+\b\lambda_i\lambda_j=0,


\end{equation}


где $\b$ --- некоторый инвариант.





После дифференцирования (8) по $x^k$ убедимся, что справедлива формула


$$


2B^2\lambda_k g_{ij}+\lambda_i d_{jk}+\lambda_j d_{ik}=0,


$$


где $d_{ij}$ --- некоторый тензор.





Из последнего равенства в случае, когда $B\neq0$, вытекает, что


$\rang\| g_{ij}\|\mr2$. Следовательно, $B=0$.


\end{pf}





Аналогичным образом можно убедиться в справедливости следующей теоремы.





\begin{thm}


Если в римановом пространстве $V_n(B)$  среди


ненулевых векторов $\lambda_i$, удовлетворяющих уравнению $(1)$, 


есть взаимоортогональные, то $B=0$.


\end{thm}








\section{Геодезические отображения четырехмерных пространств


Эйнштейна} 





Докажем теорему, которая обобщает результат А.З.\,Петрова


и Г.И.\,Голикова, сформулированный во введении.





\begin{thm}


Четырехмерные пространства Эйнштейна $V_4$, отличные от пространств


постоянной кривизны,  не допускают нетривиальные геодезические


отображения на римановы пространства~$\overline V_4$.


\end{thm}





{\bf Доказательство} проведем методом от противного.


Допустим, что четырехмерное пространство Эйнштейна $V_4$ с


непостоянной кривизной  допускает нетривиальное геодезическое


отображение на риманово пространство $\overline V_4$.





Тогда согласно известным результатам [4], [12] $\overline V_4$


является, по необходимости, также пространством Эйнштейна, а


само $V_4$ --- пространством $V_4(B)$, где $B=\frac{R}{12}$. 


В этом случае основные уравнения геодезических отображений 


запишутся в виде


\begin{equation}


\text{(а)}\ a_{ij,k} = \lambda_{i\,}g_{jk}+\lambda_{j\,}g_{ik},


\quad


\text{(б)} \ \lambda_{i,j} = \mu g_{ij}+\frac{R}{12} a_{ij},


\quad


\text{(в)} \ \mu_{,i}=\frac{R}{6}\lambda_i.


\end{equation}





Условия интегрируемости уравнений (9(б)) имеют вид


\begin{equation}


\lambda_\alpha Y_{ijk}^\alpha =0,


\end{equation}


где


$Y_{ijk}^h\equiv R_{ijk}^h-\frac R{n(n-1)}(\delta_k^hg_{ij}-


\delta_j^hg_{ik})$. 


Здесь $R_{ijk}^h$ --- тензор Римана, $\delta_i^h$ --- символы


Кронекера. 





Тензор $Y_{ijk}^h$ называется тензором {\it конциркулярной кривизны


Яно} риманова пространства $V_n$ [19].


Так как пространство $V_4$ имеет непостоянную кривизну, то его


тензор конциркулярной кривизны не равен нулю.





На основании [20]


из условий (10) при $n=4$ следует изотропность вектора 


$\lambda ^h$. Тогда из теоремы 2 вытекает $B=0$. 


В итоге скалярная кривизна $R$ равна нулю, а значит, исследуемое


пространство $V_4$ является Риччи плоским, т.\,е. имеет место


\begin{equation}


R_{ij}=0.


\end{equation}





Поскольку изотропность вектора $\lambda_i$ влечет равенство нулю


инварианта $\mu $, то


вектор $\lambda_i$ ковариантно постоянен.


В результате условия (10) принимают вид


\begin{equation}


\lambda_\alpha R_{ijk}^\alpha =0.


\end{equation}





Известно [1], [2], что в $V_4$, в котором существует


изотропный ковариантно постоянный вектор $\lambda ^h$, можно


выбрать специальную систему координат, в которой


$$


\lambda ^h=\delta_1^h,


\quad


g_{ij}(x)=


\begin{pmatrix}


0 & 0 & 0 & 1 \\


0 & g_{22} & g_{23} & g_{24} \\


0 & g_{23} & g_{33} & g_{34} \\


1 & g_{24} & g_{34} & g_{44}


\end{pmatrix}


, \quad


g^{ij}(x)=


\begin{pmatrix}


g^{11} & g^{12} & g^{13} & 1 \\


g^{12} & g^{22} & g^{23} & 0 \\


g^{13} & g^{23} & g^{33} & 0 \\


1 & 0 & 0 & 0


\end{pmatrix}. 


$$ 


 


Дальнейшие рассуждения будем вести в этой системе координат.


Тогда условия (12) принимают вид $R_{1ijk}=0$.


Затем из (11) при $i=2$, $j=4$ и $i=3$, $j=4$ имеем


\begin{align*}


g^{32}R_{3242}+g^{33}R_{3243}&=0,\\


g^{22}R_{2342}+g^{23}R_{2343}&=0.


\end{align*}


Так как


$\left| 


\begin{smallmatrix}


g^{22} & g^{23} \\ g^{23} & g^{33} \end{smallmatrix}


\right| \neq 0$,


то из последнего следует $R_{3242}=R_{3243}=0$.


Кроме того, из (11) при $i=j=2$, $i=j=3$ и $i=2$, $j=3$ будем


иметь 


$$


g^{33}R_{3223}=g^{22}R_{2332}=g^{23}R_{3232}=0.


$$


Из последнего вытекает $R_{2332}=0$.





В итоге отличными от нуля компонентами тензора


Римана (с точностью до известных тождеств тензора Римана) могут быть только


$R_{2442}$, $R_{2443}$, $R_{3443}$.





Таким образом, тензор Римана четырехмерного пространства


Эйнштейна может быть представлен в виде


\begin{equation}


R_{hijk}=\varepsilon (\xi_h\lambda_i-\xi_i\lambda_h)(\xi_j\lambda_k-\xi


_k\lambda_j).


\end{equation}


Здесь $\lambda_i\equiv\lambda ^\alpha g_{\alpha i}=


g_{1i}=\delta_i^4$, $\xi_i$ --- некоторый неколлинеарный к 


$\lambda_i$ вектор, $\varepsilon=\pm1$.





Учитывая (13), из (11) получим


\begin{equation}


\xi ^\alpha \xi_\alpha =0 \text{\ \ и \ \ }


\xi_\alpha \lambda ^\alpha =0.


\end{equation}





Ковариантно продифференцируем (13) в направлении $x^l$. В


силу ковариантного постоянства вектора $\lambda_i$ имеем


\begin{equation}


R_{hijk,l}=\varepsilon (\xi_{h,l}\lambda_i-\xi_{i,l}\lambda_h)(\xi


_j\lambda_k-\xi_k\lambda_j)+


\varepsilon (\xi_h\lambda_i-\xi_i\lambda_h)(\xi_{j,l}\lambda_k-\xi


_{k,l}\lambda_j).


\end{equation}


Учитывая тождества Бианки, проциклируем (15) по индексам 


$j$, $k$, $l$:


\begin{multline*}


(\lambda_i\xi_{h,l}-\xi_{i,l}\lambda_h)(\xi_j\lambda_k-\xi_k\lambda_j)+


(\lambda_i\xi_{h,j}-\xi_{i,j}\lambda_h)(\xi_k\lambda_l-\xi_l\lambda


_k)+ \\


+(\lambda_i\xi_{h,k}-\xi_{i,k}\lambda_h)(\xi_l\lambda_j-\xi_j\lambda_l)+


(\xi_h\lambda_i-\xi_i\lambda_h)\times \\


\times (\xi_{j,l}\lambda_k-\xi_{k,l}\lambda_j+\xi_{k,j}\lambda_l-\xi_{l,j}


\lambda_k+\xi_{l,k}\lambda_j-\xi_{j,k}\lambda_l)=0.


\end{multline*}


Из этого соотношения видно


\begin{equation}


\xi_{h,l}=\xi_h c_l+\lambda_h d_l+l_h\xi_l+f_h\lambda_l,


\end{equation}


где $c_i$, $d_i$, $l_i$, $f_i$ --- некоторые векторы.





Продифференцировав (14), учитывая (16), легко получить


\begin{equation}


\lambda ^\alpha l_\alpha =0,\quad \xi ^\alpha l_\alpha =0,\quad 


\lambda ^\alpha f_\alpha =0,\quad \xi ^\alpha f_\alpha =0.


\end{equation}





Допустим, что вектор $l_i$ линейно не выражается через векторы


$\lambda_i$ и $\xi_i$, тогда в некоторой точке $x_0$ можно


выбрать систему координат так, чтобы


$$


\lambda^i=\delta_1^i,\quad


\xi ^i=\delta_2^i,\quad l^i=\delta_3^i.


$$


Но в этом случае в силу (14), (17) и изотропности вектора


$\lambda_i$ метрика вырождается. Полученное противоречие


означает, что вектор $l_i$ можно линейно выразить через векторы


$\lambda_i$ и $\xi_i$. Аналогично установим, что и вектор $f_i$


можно линейно выразить через векторы $\lambda_i$ и $\xi_i$.


В таком случае (16) примет более простой вид


$\xi_{h,l}=\xi_hc_l+\lambda_hd_l$.


Подставив последнее в (15) и учитывая (13), получим


$$


R_{hijk,l}=\varphi_lR_{hijk},


$$


здесь $\varphi_l=2 c_l$.





Последним условием характеризуются рекуррентные или симметрические римановы


пространства, которые, как доказно в [3],


допускают нетривиальные геодезические отображения только в


том случае, когда являются пространствами постоянной кривизны.


А это противоречит сделанному предположению.$\quad\square$





В результате выделен еще один класс римановых пространств,


однозначно определенных относительно нетривиальных геодезических


отображений, какими являются, например, симметрические,


рекуррентные, обобщенно симметрические и рекуррентные и другие


пространства (см. [4], [3], [11], [13]--[22]).





\section{Об одной гипотезе А.З.\,Петрова о геодезических отображениях


пространств Эйнштейна} 





Приведем контрпример к гипотезе А.З.\,Петрова (см. введение)


([2], сс.\,355, 461).





Пусть $V_n$ $(n>4)$ --- эквидистантное пространство Эйнштейна


непостоянной кривизны с


метрикой А.\,Фиалкова [4]


\begin{equation}


ds^2=e\,dx^{1^2}+f(x^1)\,d\widetilde{s}{}^2,


\end{equation}


где $e=\pm1$, $f\neq0$, 


$d\widetilde{s}{}^2=\widetilde{g}_{\alpha \beta


}(x^2,\dotsc,x^n)dx^\alpha dx^\beta$


$(\alpha ,\beta >1)$


--- метрика некоторого пространства $\widetilde{V}_{n-1}$,


которое по необходимости является 


пространством Эйнштейна, отличного от пространства постоянной


кривизны, и функция $f(x^1)$ удовлетворяет условию


$$


f=


\begin{cases}


\wt B/B \cos^2(\sqrt{-eB}x^1+b), & eB<0, \ \ \wt B\nr0; \\


\wt B/B (a+{\sh}^2(\sqrt{eB}x^1+b)), & eB>0, \ \ \wt B\nr0;\\


b \exp^2(\sqrt{eB}x^1), & eB\ge 0, \ \ \wt B=0; \\


-e\wt B(x^1+b)^2, &  \phantom{e}B=0, \ \ \wt B \nr0,


\end{cases}


$$


где $a=0,1$, $b$ --- $\const$, $B=\frac{R}{n(n-1)}$,


$\wt B=\frac{\wt R}{(n-1)(n-2)}$, $R$ $(\wt R)$ ---


скалярные кривизны пространств\linebreak \vn\ $(\widetilde{V}_{n-1})$. 





Доказано [4], [13], что пространство \vn, отнесенное к системе


координат (18), допускает геодезическое отображение на риманово


пространство \vnn, метрическая форма которого имеет вид


$$


d\ov s^2=\frac{ep}{(1+qf)^2}(dy^1)^2+\frac{fp}{1+qf}d\wt s^2,


$$


где $p$, $q$ --- некоторые постоянные такие, что


$p\nr 0$, $1+qf\nr 0$.


При $q f^{\prime }\not\equiv 0$ 


отображение является нетривиальным. Координаты $x$ являются


общими по этому отображению. 





Сигнатуры метрик \vn\ и \vnn\ различны, когда $1+qf<0$, в


противном случае совпадают.





Легко видеть, что при подходящем подборе постоянных $e$, $q$ можно


построить пример нетривиального геодезического отображения между


эйнштейновыми пространствами непостоянной кривизны с сигнатурой


Минковского и размерности выше четырех. Это и есть контрпример к


приведенной гипотезе А.З.\,Петрова.





\begin{thebibliography}{99}





\bibitem{eis}%1


Эйзенхарт Л.П. {\em Риманова геометрия}. -- М.: Ин. лит., 1948. -- 315 с.





\bibitem{pet}%2


Петров А.З. {\em Новые методы в общей теории относительности}.


-- М.: Наука, 1966. -- 495~с. 





\bibitem{si}%3


Синюков Н.С. {\em Геодезические отображения римановых пространств}.


-- М.: Наука, 1979. --~256~с.





\bibitem{mige}%4


Mike\v s J. {\em Geodesic mappings of affine-connected and Riemannian spaces}


// J.~Math. Sci. New York. -- 1996. -- V.\,78. -- \No\,2. -- P.\,311--333.





\bibitem{pet49}%5


Петров А.З. {\em О геодезических отображениях римановых пространств


неопределенной метрики}. -- Учен. зап. Казанск. ун-та. --


1949. -- T.\,109. -- C.\,4--36.





\bibitem{pet61}%6


Петров А.З. {\em О геодезическом отображении  пространств Эйнштейна} //


Изв. вузов. Математика. -- 1961. -- \No\,2. -- C.\,130--136.





\bibitem{gold}%7 


Голиков В.И. {\em Поля тяготения с общими геодезическими}:


Дис. $\dotsc$ канд.


физ.-матем. наук. -- Казанск. ун-т, Казань, 1963.





\bibitem{gol63}%8 


Голиков В.И. {\em Поля тяготения с общими геодезическими}. I //


Учен. зап. Казанск. ун-та. 


-- 1963. -- T.\,2. -- C.\,72--95.





\bibitem{9}


Голиков В.И. {\em Поля тяготения с общими геодезическими}. II //


Учен. зап. Казанск. ун-та. 


-- 1963. -- T.\,12. -- C.\,59--67.





\bibitem{venei}%9


Venzi P. {\em On geodesic mappings in Riemannian and pseudo-Riemannian


manifolds} // Tensor. -- 1979. -- V.\,33. -- P.\,23--28.





\bibitem{mid}%10


Микеш Й. {\em Геодезические и голоморфно-проективные отображения


специальных римановых пространств}: Дис.~$\dotsc$ канд.


физ.-матем. наук. -- Одесск. ун-т, 1979. -- 107 с.





\bibitem{m11}%11


Микеш Й. {\em О геодезических отображениях пространств Эйнштейна} //


Матем. заметки. -- 1980. -- T.\,28. -- \No\,6. -- C.\,935--938.





\bibitem{midd}%12


Mike\v s J. {\em Geodetick\'a, $F$-plan\'arn\'\i\ a holomorfn\v e


projektivn\'\i\ zobrazen\'\i\ Riemannov\'ych variet a variet s


afinn\'\i\ konex\'\i}\/: Dokt. disert., Olomouc, 1995. -- 181 p.





\bibitem{sol56}%14


Солодовников А.С. {\em Геодезические классы пространств $V(K)$}


// ДАН СССР. -- 1956. -- T.\,111, -- \No\,1. -- C.\,33--36.





\bibitem{K18}%13


Микеш Й., Киосак В.А. {\em О геодезических отображениях специальных


пространств}. -- Киев, 1985. -- 24 с. -- Деп. в УкрНИИТИ


05.05.1985, \No\,904-24Ук.





\bibitem{kru67}%15.


Кручкович Г.И. {\em О пространствах $V(K)$


и их геодезических отображениях} //


Тр. Всесоюзн. заочн. энерг. ин-та.  -- М., 1967. -- T.\,33. -- C.\,3--18.





\bibitem{mikio}%16


Микеш Й., Киосак В.А. {\em О степени подвижности римановых пространств


относительно геодезических отображений} //


Геометрия погружен. многообразий. -- М.: МГПИ, 1986. -- C.\,35--39.





\bibitem{gor}%17


Горбатый Е.З. {\em О геодезическом отображении эквидистантных


римановых пространств и пространств первого класса} //


Укр. геометрич. сб. -- 1972. -- T.\,12. -- C.\,45--53.





\bibitem{yan}%18


Yano K. {\em The theory of Lie derivatives and its applications}. --


Amsterdam: North Holland Publ. Gr\"oningen, Nordhoff, 1957. -- 293 c.





\bibitem{shostr}%19


Схоутен И.А., Стройк Д.Дж. {\em Введение в новые методы


дифференциальной геометрии}. I. -- М.--Л.: Гостехиздат, 1939. -- 181 c.





\bibitem{21}


Схоутен И.А., Стройк Д.Дж. {\em Введение в новые методы дифференциальной


геометрии}. II. -- М.: Ин. лит., 1948. -- 348 с.





\bibitem{ramiga}%20


Радулович Ж., Микеш Й., Гаврильченко М.Л. {\em Геодезические отображения и


деформации римановых пространств}. --


Podgorica, Одесса: Изд. Одесск. ун-та, 1997. -- 127 с.





\end{thebibliography}





\vskip 0.5cm





\post{\it Одесский государственный}{\it Поступила}


\post{\it университет}{07.12.2002}


\post{\it Оломоуцкий университет}{}


\post{\it $($Чешская республика$)$}{}





\label{end}


\end{document}


